\section{Background}
\label{sec:background}

\subsection{Terminology}
Directions on a boat are relative to the helm: \emph{forward} (fore) points
toward the front and \emph{aft} toward the back, while \emph{outboard} points
away from the centerline and \emph{inboard} toward it. The front of the vessel
is the \emph{bow}, the rear is the \emph{stern}, and the right and left sides
are \emph{starboard} and \emph{port}. The \emph{hull} is the body of the boat;
the \emph{keel} is the fully submerged fin beneath it; the \emph{rudder} is the
small steering plate at the stern; and the \emph{mast} is the vertical member
that carries the sail. \emph{Freeboard} is the height of hull above the
waterline and \emph{draft} is the depth below it. A self-trimming wingsail
consists of a rigid main wing, a smaller
\emph{tail} (also called a vane) carried aft on a boom, a counterweight, and
the pivot about which the assembly rotates. Fig.~\ref{fig:anatomy} labels
these terms on ALAN's layout.

\begin{figure}[!t]
\centering
\subfloat[]{%
\begin{tikzpicture}[scale=0.8, every node/.style={font=\scriptsize}]
  % hull
  \draw[fill=gray!25] (-2.0,0) -- (-2.0,0.32) .. controls (0.2,0.42) and (1.6,0.35) .. (2.6,0)
                       .. controls (1.6,-0.35) and (0.2,-0.42) .. (-2.0,-0.32) -- cycle;
  \draw[dashed,gray] (-2.4,0) -- (3.2,0);
  \fill (0.35,0) circle (1.5pt);
  \node[anchor=north west,inner sep=1pt] at (0.4,0) {mast};
  \node[anchor=east] at (-2.05,0.25) {stern};
  \node[anchor=south west,inner sep=1pt] at (2.6,0.02) {bow};
  \node at (0.3,0.95) {port};
  \node at (0.3,-0.95) {starboard};
  \node[gray,anchor=north west,inner sep=1pt] at (2.6,-0.02) {centerline};
  % direction arrows
  \draw[{Stealth[length=1.6mm]}-{Stealth[length=1.6mm]}] (-1.6,-1.6) -- (1.6,-1.6);
  \node[anchor=east] at (-1.6,-1.6) {aft};
  \node[anchor=west] at (1.6,-1.6) {fore};
  \draw[-{Stealth[length=1.6mm]}] (3.9,0.45) -- (3.9,1.25) node[above] {outboard};
  \draw[-{Stealth[length=1.6mm]}] (3.9,-1.25) -- (3.9,-0.45);
  \node[anchor=north] at (3.9,-1.25) {inboard};
\end{tikzpicture}}

\subfloat[]{%
\begin{tikzpicture}[scale=0.8, every node/.style={font=\scriptsize}]
  % water
  \fill[blue!8] (-2.6,-2.3) rectangle (3.9,0);
  \draw[blue!60!black] (-2.6,0) -- (3.9,0) node[anchor=south east,inner sep=1pt] {waterline};
  % hull profile
  \draw[fill=gray!25] (-2.0,0.32) -- (2.5,0.32) .. controls (2.3,0.0) and (1.7,-0.32) .. (0.8,-0.36)
                       -- (-1.9,-0.3) -- cycle;
  % keel and bulb
  \draw[fill=gray!45] (0.1,-0.34) -- (0.55,-0.36) -- (0.45,-1.7) -- (0.1,-1.7) -- cycle;
  \draw[fill=gray!60] (-0.35,-1.62) .. controls (-0.2,-1.45) and (0.9,-1.45) .. (1.05,-1.72)
                       .. controls (0.9,-1.95) and (-0.2,-1.95) .. (-0.35,-1.72) -- cycle;
  \node[anchor=west] at (0.6,-1.15) {keel and bulb};
  % rudder
  \draw[fill=gray!45] (-1.75,-0.3) -- (-1.45,-0.3) -- (-1.5,-1.0) -- (-1.8,-1.0) -- cycle;
  \node[anchor=east] at (-1.85,-0.75) {rudder};
  % mast
  \draw[thick] (0.35,0.32) -- (0.35,1.7);
  \node[anchor=west] at (0.4,1.5) {mast};
  % freeboard and draft
  \draw[{Stealth[length=1.2mm]}-{Stealth[length=1.2mm]}] (-2.2,0) -- (-2.2,0.32);
  \node[anchor=east,inner sep=1pt] at (-2.25,0.2) {freeboard};
  \draw[{Stealth[length=1.4mm]}-{Stealth[length=1.4mm]}] (3.3,0) -- (3.3,-1.95);
  \draw[densely dotted] (-0.35,-1.95) -- (3.4,-1.95);
  \node[anchor=west] at (3.35,-1.0) {draft};
  \node[anchor=south] at (-1.6,0.32) {stern};
  \node[anchor=south] at (2.1,0.32) {bow};
\end{tikzpicture}}
\caption{Vessel terminology on ALAN's layout (schematic, not to scale).
(a) Top view of the hull with the fore--aft and inboard--outboard
directions. Port is on the left when facing the bow.
(b) Side view: freeboard and draft are measured from the waterline; the fin
keel and bulb sit below the mast, and the spade rudder at the stern.}
\label{fig:anatomy}
\end{figure}

\subsection{Physics of Sailing}
A sailboat is propelled by two opposing lift forces~\cite{skene2001}. Wind
flowing past the sail produces lift perpendicular to the apparent wind, and
on most headings this force points largely sideways. The keel, moving
through the water at a small leeway angle, produces a second lift force that
opposes the sideways component. What remains is a forward drive force, as
shown in Fig.~\ref{fig:physics}(b).

Because of this two-force balance, both the \emph{point of sail} (the boat's
heading relative to the wind) and the sail's \emph{angle of attack} (the angle
between the wing chord and the apparent wind) determine boat speed. A boat
cannot make headway within roughly \SI{45}{\degree} of the true wind on either
side; this \SI{90}{\degree} sector is called \emph{in irons} or the
\emph{no-go zone} (Fig.~\ref{fig:physics}(a)). To reach a target inside the
no-go zone, the boat must \emph{tack}, sailing a zig-zag course at the edge of
the zone.

The rudder acts as a second, smaller foil at the stern. Turning it toward
the wind reduces the effective keel force aft and turns the bow away from the
wind; turning it away from the wind does the opposite. The point at which the
sail force acts is the \emph{center of effort} (COE), and the point at which the
keel force acts is the \emph{center of lateral resistance} (CLR). Placing the
COE slightly aft of the CLR gives a mild \emph{weather helm}: the boat tends to
turn into the wind and depower if left uncontrolled, which is a safe default
for an autonomous vessel.

\begin{figure}[!t]
\centering
\subfloat[]{%
\begin{tikzpicture}[scale=0.72, every node/.style={font=\scriptsize}]
  % no-go zone
  \fill[red!15] (0,0) -- (45:2.5) arc (45:135:2.5) -- cycle;
  \node at (90:2.1) {no-go zone};
  % wind
  \foreach \x in {-0.6,0,0.6}
    \draw[-{Latex[length=2mm]},thick,blue!70!black] (\x,3.6) -- (\x,2.85);
  \node[blue!70!black,anchor=east] at (-0.8,3.25) {true wind};
  % headings
  \foreach \a/\lab/\anc in {45/close-hauled/west, 135/close-hauled/east,
                            0/beam reach/north, 180/beam reach/north,
                            -45/broad reach/north, -135/broad reach/north,
                            -90/run/north}
    \draw[-{Stealth[length=2mm]}] (0,0) -- (\a:2.2);
  \node[anchor=south west] at (45:2.2) {close-hauled};
  \node[anchor=south east] at (135:2.2) {close-hauled};
  \node[anchor=north] at (0:2.0) {\shortstack{beam\\reach}};
  \node[anchor=north] at (180:2.0) {\shortstack{beam\\reach}};
  \node[anchor=north west] at (-45:2.1) {broad reach};
  \node[anchor=north east] at (-135:2.1) {broad reach};
  \node[anchor=north] at (-90:2.2) {run};
  \draw[dashed] (0,0) -- (90:2.5);
  \draw (90:1.1) arc (90:45:1.1); \node at (67.5:1.45) {$45^\circ$};
  \fill (0,0) circle (1.5pt);
\end{tikzpicture}}
\hfil
\subfloat[]{%
\begin{tikzpicture}[scale=0.72, every node/.style={font=\scriptsize}]
  % hull, heading +x
  \draw[fill=gray!20] (-1.6,0) .. controls (-1.2,0.45) and (0.9,0.45) .. (1.7,0)
                       .. controls (0.9,-0.45) and (-1.2,-0.45) .. (-1.6,0);
  \foreach \x in {-0.6,0,0.6}
    \draw[-{Latex[length=2mm]},thick,blue!70!black] (\x+0.3,2.9) -- (\x+0.3,2.2);
  \node[blue!70!black] at (1.75,2.55) {wind};
  % sail force (lift)
  \coordinate (C) at (0.1,0);
  \draw[-{Stealth[length=2.5mm]},very thick,red!70!black] (C) -- ++(-62:2.3) node[below] {$F_S$};
  % components
  \draw[-{Stealth[length=2mm]},dashed,red!70!black] (C) -- ++(1.08,0) node[above right=-2pt] {$F_D$};
  \draw[-{Stealth[length=2mm]},dashed,red!70!black] (C) -- ++(0,-2.03) node[left] {$F_H$};
  % keel force
  \draw[-{Stealth[length=2.5mm]},very thick,green!50!black] (C) -- ++(0,2.03) node[left] {$F_K$};
  \draw[-{Stealth[length=2mm]},thick] (1.9,0) -- (2.8,0) node[right] {motion};
\end{tikzpicture}}
\caption{Sailing fundamentals. (a) Points of sail; headings within
\SI{45}{\degree} of the true wind are not sailable. (b) Force balance seen from
above: the sail force $F_S$ splits into a drive component $F_D$ and a heeling
component $F_H$, which the keel force $F_K$ cancels.}
\label{fig:physics}
\end{figure}

\subsection{Hull, Keel, and Rudder Options}
Hulls are broadly \emph{displacement} or \emph{planing}. Displacement hulls
are rounder and deeper and float by displacing water; they are slower but
more stable. Planing hulls are flatter and rise onto the surface at speed;
they are faster but less stable~\cite{skene2001}. Keels are either
\emph{full}, running the length of the hull with shallow draft and high drag,
or \emph{fin}, a thin deep plate that is more efficient but increases
draft~\cite{sailingcats}. Rudders are classified by balance: a fully balanced
rudder pivots about its center of pressure and is easy to turn but does not
self-center, while an unbalanced rudder pivots at its leading edge, needs more
torque, and returns to center on its own~\cite{noahrudder}. A \emph{spade}
rudder hangs from its stock alone; a \emph{skeg} rudder is additionally
supported by a fin bolted to the hull.

\subsection{Wingsail Control Strategies}
Cloth sails are inexpensive but wear quickly. Rigid sails made of foam,
plastic, or composite are more durable and generate lift more efficiently,
which is why most autonomous sailboats use them~\cite{silva2019,sun2025}.
Three control strategies exist for rigid wingsails~\cite{sun2025}, compared
in Table~\ref{tab:sailtypes}. A \emph{traditional} wingsail is rotated
directly by a motor at the mast, which is simple but requires enough torque
to turn and hold the whole wing. A \emph{flap} wingsail adds a hinged trailing
flap; the motor moves only the flap, and the aerodynamic force on the flap
rotates the wing. A \emph{self-trimming} wingsail carries a tail roughly one
third the size of the main wing on a boom aft of the pivot. The motor turns
only the tail, whose aerodynamic moment rotates the wing to an equilibrium
angle of attack~\cite{elkaim2008,sailwings}. Actuation power is lowest for
the self-trimming design, at the cost of an indirect, delayed response that is
harder to control.

\begin{table}[!t]
\caption{Comparison of Rigid Wingsail Control Strategies}
\label{tab:sailtypes}
\centering
\begin{tabular}{@{}lccc@{}}
\toprule
 & Traditional & Flap & Self-trimming \\
\midrule
Actuated surface    & Whole wing & Flap & Tail \\
Actuator power      & High       & Medium & Low \\
Control complexity  & Low        & Medium & High \\
Response to command & Direct     & Indirect & Indirect, delayed \\
\bottomrule
\end{tabular}
\end{table}
